\subsection{Eigenfunction Expansions}\label{sec:elliptic-eigenfunction-expansions}

Under suitable hypotheses on the elliptic operator $L$, Theorem~\ref{thm:weak-dirichlet-positive-shift-existence} guarantees that there exists $\lambda_G$ such that if $\widetilde\lambda>\lambda_G$, then for any $f\in H^{-k}(\Omega)$ there exists a unique (weak) solution $u\in H_0^k(\Omega)$ of the Dirichlet problem for\[\widetilde L(\bm{x},D)u:=L(\bm{x},D)u+\widetilde\lambda u=f.\]In this section we will apply some of the operator techniques developed in the previous chapter to this problem. This investigation will give us two basic improvements over the present existence theory. First, the Fredholm theorems will give us information on the existence and uniqueness of solutions for values of $\widetilde\lambda<\lambda_G$. Second, if the operator $L$ satisfies a symmetry condition, we can use the method of eigenfunction expansion to construct (or in real life approximate) solutions.

\subsubsection{Fredholm Theory}\label{subsec:elliptic-fredholm-theory}

In this section we consider the nonhomogeneous eigenvalue problem
\begin{equation}\label{eq:elliptic-nonhomogeneous-shifted-eigenvalue-problem}
L(\bm{x},D)u+\lambda u=f
\end{equation}
for $f\in L^2(\Omega)$, where $L(\bm{x},D)$ is the operator\[L(\bm{x},D)u=\sum_{0\leq|\sigma|,|\gamma|\leq k}(-1)^{|\sigma|}D^\sigma(a_{\sigma\gamma}(\bm{x})D^\gamma u),\]and the bilinear form associated with $L$ is\[B[v,u]=\sum_{0\leq|\sigma|,|\gamma|\leq k}\int_\Omega a_{\sigma\gamma}(\bm{x})D^\gamma uD^\sigma v\,\mathrm{d}\bm{x}.\]

Let us assume the hypotheses of Theorem~\ref{thm:weak-dirichlet-positive-shift-existence} are satisfied and fix $\widetilde\lambda>\lambda_G$ with $\widetilde\lambda>0$. Then for any $f\in L^2(\Omega)$ there is a unique solution $u\in H_0^k(\Omega)$ to the problem
\begin{equation}\label{eq:elliptic-shifted-solution-variational-definition}
B_{\widetilde\lambda}[v,u]:=B[v,u]+\widetilde\lambda(v,u)=(v,f)_{L^2(\Omega)}\quad\text{for every }v\in H_0^k(\Omega).
\end{equation}
We now define an operator $\overline G:L^2(\Omega)\to H_0^k(\Omega)$ as follows: for every $f\in L^2(\Omega)$ we define
\begin{equation}\label{eq:elliptic-sobolev-solution-operator-definition}
\overline G(f):=\widetilde\lambda u,
\end{equation}
where $u$ is the unique (weak) solution of the Dirichlet problem for
\begin{equation}\label{eq:elliptic-solution-operator-shifted-dirichlet-equation}
L(\bm{x},D)u+\widetilde\lambda u=f;
\end{equation}
i.e., $u$ solves \eqref{eq:elliptic-shifted-solution-variational-definition}. In other words, for every $f\in L^2(\Omega)$ and $v\in H_0^k(\Omega)$ we have
\begin{equation}\label{eq:elliptic-sobolev-solution-operator-variational-identity}
B_{\widetilde\lambda}[v,\overline G(f)]=\widetilde\lambda(v,f)_{L^2(\Omega)}.
\end{equation}


Formally, we have
\begin{equation}\label{eq:elliptic-sobolev-solution-operator-formal-inverse}
\overline G=\widetilde\lambda(L+\widetilde\lambda)^{-1}.
\end{equation}
By \eqref{eq:weak-dirichlet-solution-negative-norm-bound}, the operator $\overline G$ is bounded. We now define the operator $G:L^2(\Omega)\to L^2(\Omega)$ by the composition of $\overline G$ and $\overline I$,
\begin{equation}\label{eq:elliptic-l2-compact-solution-operator-definition}
G:=\overline I\,\overline G,
\end{equation}
where $\overline I$ is the identity mapping from $H_0^k(\Omega)$ to $L^2(\Omega)$. We know from Theorem~\ref{thm:compact-successive-sobolev-orders} that this operator is compact. Since the composition of a bounded operator and a compact operator is compact (cf. Problem~\ref{ex:bounded-compact-operator-products}) we have the following.

\begin{Lemma}
\label{lem:elliptic-l2-solution-operator-compact}
The solution operator $G:L^2(\Omega)\to L^2(\Omega)$ is compact.
\end{Lemma}
